\appendix
\section*{Appendix}
\addcontentsline{toc}{section}{Appendix}

\section{System Prompts and Agent Directives}
\label{app:prompts}

This appendix presents the exact system directives and prompt templates employed by the agent modules.

\subsection{Domain Classification Directives}

\begin{lstlisting}[language=bash, caption={Domain Classification System Directive.}, label={lst:prompt_domain}]
System Directive:
You are an expert medical domain classifier. Analyze the provided clinical vignette and identify up to 5 primary medical subspecialties relevant to diagnosing the case (e.g., Cardiology, Nephrology, Pathology, Pharmacology, Neurology).

Format your output strictly as a comma-separated list of domains:
Domains: [Domain 1], [Domain 2], [Domain 3]
\end{lstlisting}

\subsection{Consensus Verification Directives}

\begin{lstlisting}[language=bash, caption={Verifier Agent Voting Directive.}, label={lst:prompt_verifier}]
System Directive:
You are a senior clinical reviewer representing the domain specialist: {domain_name}.
Review the following diagnostic synthesis report for the clinical vignette.

Determine if the diagnosis and option selection are factually accurate and logically sound from your specialty perspective.

Output your response strictly in the format:
Vote: [YES or NO]
Advice: [Provide specific revision instructions if NO, otherwise write NONE]
\end{lstlisting}

\section{Data Schemas and Record Specifications}
\label{app:schemas}

Each prediction record saved to \texttt{outputs/<VARIANT>\_predictions.jsonl} follows the JSON schema specification detailed in Listing~\ref{lst:json_schema}.

\begin{lstlisting}[language=Python, caption={Serialized JSON Prediction Record Schema.}, label={lst:json_schema}]
{
  "$schema": "http://json-schema.org/draft-07/schema#",
  "title": "MedRAGAgentsPredictionRecord",
  "type": "object",
  "properties": {
    "question": { "type": "string" },
    "options": { "type": "object" },
    "gold_answer": { "type": "string" },
    "pred_answer": { "type": "string" },
    "raw_output": { "type": "string" },
    "snippets": { "type": "array" },
    "scores": { "type": "array" },
    "question_domains": { "type": "array" },
    "option_domains": { "type": "array" },
    "syn_report": { "type": "string" },
    "vote_history": { "type": "array" },
    "revision_history": { "type": "array" },
    "from_cache": { "type": "boolean" }
  },
  "required": ["question", "options", "gold_answer", "pred_answer"]
}
\end{lstlisting}

\section{Derivation of Non-Parametric Bootstrap Algorithm}
\label{app:bootstrap}

Listing~\ref{lst:bootstrap_code} details the non-parametric bootstrap estimation procedure used to derive 95\% confidence intervals for classification accuracy.

\begin{lstlisting}[language=Python, caption={Non-Parametric Bootstrap Resampling Algorithm in evaluate.py.}, label={lst:bootstrap_code}]
import numpy as np

def bootstrap_confidence_interval(y_true, y_pred, n_bootstraps=1000, ci=95):
    """Computes non-parametric bootstrap confidence interval for accuracy."""
    n_samples = len(y_true)
    bootstrapped_accuracies = []
    rng = np.random.RandomState(42)

    for _ in range(n_bootstraps):
        indices = rng.choice(n_samples, size=n_samples, replace=True)
        sample_true = [y_true[i] for i in indices]
        sample_pred = [y_pred[i] for i in indices]
        acc = sum(t == p for t, p in zip(sample_true, sample_pred)) / n_samples
        bootstrapped_accuracies.append(acc)

    lower_p = (100 - ci) / 2.0
    upper_p = 100 - lower_p
    lower = np.percentile(bootstrapped_accuracies, lower_p)
    upper = np.percentile(bootstrapped_accuracies, upper_p)
    return lower, upper
\end{lstlisting}
